\subsection{7.17 State-Space Formulation and Explicit Derivation}

This subsection provides a detailed step-by-step derivation of the state-space representation shown in the attachment, transforming the Euler-Lagrange equations into explicit first-order differential equations.

\subsubsection{7.17.1 State Vector Definition}

We define the state vector as:

\begin{equation}
\boxed{\mathbf{x} = \begin{bmatrix}
x_m \\
y_m \\
x_{m,1} \\
y_{m,1} \\
T_1 \\
T_2
\end{bmatrix} = \begin{bmatrix}
x_m \\
y_m \\
\dot{x}_m \\
\dot{y}_m \\
T_1 \\
T_2
\end{bmatrix}}
\end{equation}

where:
\begin{itemize}
    \item $(x_m, y_m)$ = End-effector (cable attachment point M) position
    \item $(\dot{x}_m, \dot{y}_m)$ = End-effector velocity (denoted as $x_{m,1}, y_{m,1}$ in the notation of the attachment)
    \item $T_1, T_2$ = Cable tensions
\end{itemize}

\textbf{Alternative notation mapping:}
\begin{align}
x_m &\leftrightarrow x_m \text{ (position)} \\
y_m &\leftrightarrow y_m \text{ (position)} \\
x_{m,1} &\leftrightarrow \dot{x}_m \text{ (velocity)} \\
y_{m,1} &\leftrightarrow \dot{y}_m \text{ (velocity)} \\
T_1, T_2 &\leftrightarrow T_1, T_2 \text{ (cable tensions)}
\end{align}

\subsubsection{7.17.2 Forward Kinematics}

The forward kinematics relate joint angles to end-effector position. For a 2-link serial arm with cable attachment at the tip M:

\begin{equation}
\begin{bmatrix}
x_m \\
y_m
\end{bmatrix} = \begin{bmatrix}
l_1 \cos(q_1) + l_2 \cos(q_1 + q_2) \\
l_1 \sin(q_1) + l_2 \sin(q_1 + q_2)
\end{bmatrix}
\end{equation}

Taking time derivatives (velocity kinematics):

\begin{equation}
\begin{bmatrix}
\dot{x}_m \\
\dot{y}_m
\end{bmatrix} = \mathbf{J}_q(\mathbf{q}) \dot{\mathbf{q}}
\end{equation}

where the Joint Jacobian is:

\begin{equation}
\mathbf{J}_q(\mathbf{q}) = \begin{bmatrix}
-l_1 \sin(q_1) - l_2 \sin(q_1+q_2) & -l_2 \sin(q_1+q_2) \\
l_1 \cos(q_1) + l_2 \cos(q_1+q_2) & l_2 \cos(q_1+q_2)
\end{bmatrix}
\end{equation}

Therefore:

\begin{align}
\dot{x}_m &= -l_1 \sin(q_1) \dot{q}_1 - l_2 \sin(q_1+q_2) (\dot{q}_1 + \dot{q}_2) \\
\dot{y}_m &= l_1 \cos(q_1) \dot{q}_1 + l_2 \cos(q_1+q_2) (\dot{q}_1 + \dot{q}_2)
\end{align}

\textbf{State equations (1) and (2):}

\begin{equation}
\boxed{\begin{aligned}
\dot{x}_m &= -l_1 \sin(q_1) \dot{q}_1 - l_2 \sin(q_1+q_2) (\dot{q}_1 + \dot{q}_2) \\
\dot{y}_m &= l_1 \cos(q_1) \dot{q}_1 + l_2 \cos(q_1+q_2) (\dot{q}_1 + \dot{q}_2)
\end{aligned}}
\end{equation}

These are the first two equations in the attachment's state vector.

\subsubsection{7.17.3 Acceleration Kinematics}

Taking time derivatives of velocity kinematics:

\begin{equation}
\ddot{\mathbf{r}}_m = \mathbf{J}_q \ddot{\mathbf{q}} + \dot{\mathbf{J}}_q \dot{\mathbf{q}}
\end{equation}

where $\dot{\mathbf{J}}_q$ is the time derivative of the Jacobian:

\begin{equation}
\dot{\mathbf{J}}_q = \begin{bmatrix}
-l_1 \cos(q_1) \dot{q}_1 - l_2 \cos(q_1+q_2)(\dot{q}_1 + \dot{q}_2) & -l_2 \cos(q_1+q_2)(\dot{q}_1 + \dot{q}_2) \\
-l_1 \sin(q_1) \dot{q}_1 - l_2 \sin(q_1+q_2)(\dot{q}_1 + \dot{q}_2) & -l_2 \sin(q_1+q_2)(\dot{q}_1 + \dot{q}_2)
\end{bmatrix}
\end{equation}

Therefore:

\begin{equation}
\begin{bmatrix}
\ddot{x}_m \\
\ddot{y}_m
\end{bmatrix} = \mathbf{J}_q \begin{bmatrix}
\ddot{q}_1 \\
\ddot{q}_2
\end{bmatrix} + \dot{\mathbf{J}}_q \begin{bmatrix}
\dot{q}_1 \\
\dot{q}_2
\end{bmatrix}
\end{equation}

\textbf{Expanding the acceleration components:}

\begin{equation}
\ddot{x}_m = -l_1 \sin(q_1) \ddot{q}_1 - l_2 \sin(q_1+q_2) (\ddot{q}_1 + \ddot{q}_2) + \text{nonlinear terms}
\end{equation}

\begin{equation}
\ddot{y}_m = l_1 \cos(q_1) \ddot{q}_1 + l_2 \cos(q_1+q_2) (\ddot{q}_1 + \ddot{q}_2) + \text{nonlinear terms}
\end{equation}

\textbf{State equations (3) and (4):}

\begin{equation}
\boxed{\begin{aligned}
\dot{x}_{m,1} &= \ddot{x}_m \\
\dot{y}_{m,1} &= \ddot{y}_m
\end{aligned}}
\end{equation}

These define the time derivatives of the velocity components, which equal accelerations.

\subsubsection{7.17.4 Dynamics Equations from Euler-Lagrange}

From Section 7.6, the Euler-Lagrange equations yield:

\begin{equation}
\mathbf{M}(\mathbf{q}) \ddot{\mathbf{q}} + \mathbf{G}(\mathbf{q}) = \boldsymbol{\tau} + \mathbf{J}_{\rho}^T(\mathbf{q}) \mathbf{T}
\end{equation}

Solving for accelerations:

\begin{equation}
\ddot{\mathbf{q}} = \mathbf{M}^{-1} \left[ \boldsymbol{\tau} + \mathbf{J}_{\rho}^T(\mathbf{q}) \mathbf{T} - \mathbf{G}(\mathbf{q}) \right]
\end{equation}

For our planar 2-DOF arm with constant mass matrix $\mathbf{M}$:

\begin{equation}
\begin{bmatrix}
\ddot{q}_1 \\
\ddot{q}_2
\end{bmatrix} = \mathbf{M}^{-1} \begin{bmatrix}
\tau_1 + J_{\rho 1,1} T_1 + J_{\rho 1,2} T_2 - g_1 \\
\tau_2 + J_{\rho 2,1} T_1 + J_{\rho 2,2} T_2 - g_2
\end{bmatrix}
\end{equation}

\textbf{Explicit form for 2-DOF case:}

Let $m_1$ and $m_2$ be the link masses, and compute:

\begin{equation}
\mathbf{M}^{-1} = \frac{1}{\det(\mathbf{M})} \begin{bmatrix}
I_2 & -I_2 \\
-I_2 & I_1 + I_2
\end{bmatrix}
\end{equation}

where $\det(\mathbf{M}) = I_1 I_2 - I_2^2 = I_2(I_1 - I_2)$.

This gives:

\begin{equation}
\ddot{q}_1 = \frac{1}{I_1 I_2 - I_2^2} \left[ I_2(\tau_1 + \mathbf{J}_{\rho,1}^T \mathbf{T} - g_1) - I_2(\tau_2 + \mathbf{J}_{\rho,2}^T \mathbf{T} - g_2) \right]
\end{equation}

\begin{equation}
\ddot{q}_2 = \frac{1}{I_1 I_2 - I_2^2} \left[ -I_2(\tau_1 + \mathbf{J}_{\rho,1}^T \mathbf{T} - g_1) + (I_1+I_2)(\tau_2 + \mathbf{J}_{\rho,2}^T \mathbf{T} - g_2) \right]
\end{equation}

\subsubsection{7.17.5 Cable Tension Dynamics}

In the attachment, $T_1$ and $T_2$ evolve according to cable actuation dynamics. Without additional actuator models, we typically assume:

\begin{equation}
\dot{T}_i = k_i (\tau_{c,i} - T_i) + d_i \dot{\rho}_i
\end{equation}

where:
\begin{itemize}
    \item $\tau_{c,i}$ = Cable tension command (control input)
    \item $k_i$ = Tension actuator response rate
    \item $d_i$ = Damping coefficient related to cable compliance
    \item $\dot{\rho}_i$ = Rate of change of cable length (velocity along cable direction)
\end{itemize}

\textbf{For ideal tension control (fast response):}

\begin{equation}
\boxed{\begin{aligned}
\dot{T}_1 &= 0 \\
\dot{T}_2 &= 0
\end{aligned}}
\end{equation}

This assumes tensions are directly controllable (quasi-static or ideal controller).

\textbf{For realistic tension dynamics:}

\begin{equation}
\boxed{\begin{aligned}
\dot{T}_1 &= -k_1 T_1 + k_1 \tau_{c,1} \\
\dot{T}_2 &= -k_2 T_2 + k_2 \tau_{c,2}
\end{aligned}}
\end{equation}

with time constants $\tau_c = 1/k_i$.

\textbf{State equations (5) and (6) from attachment:}

\begin{equation}
\boxed{\begin{aligned}
\dot{T}_1 &= 0 \text{ (or first-order dynamics as above)} \\
\dot{T}_2 &= 0 \text{ (or first-order dynamics as above)}
\end{aligned}}
\end{equation}

\subsubsection{7.17.6 Complete State-Space System}

The complete state-space representation is:

\begin{equation}
\dot{\mathbf{x}} = \mathbf{f}(\mathbf{x}, \mathbf{u})
\end{equation}

\textbf{Explicitly:}

\begin{equation}
\boxed{\begin{aligned}
\dot{x}_m &= -l_1 \sin(q_1) \dot{q}_1 - l_2 \sin(q_1+q_2) (\dot{q}_1 + \dot{q}_2) \\
\dot{y}_m &= l_1 \cos(q_1) \dot{q}_1 + l_2 \cos(q_1+q_2) (\dot{q}_1 + \dot{q}_2) \\
\dot{x}_{m,1} &= \ddot{x}_m(\mathbf{q}, \dot{\mathbf{q}}, \mathbf{T}, \boldsymbol{\tau}) \\
\dot{y}_{m,1} &= \ddot{y}_m(\mathbf{q}, \dot{\mathbf{q}}, \mathbf{T}, \boldsymbol{\tau}) \\
\dot{T}_1 &= T_1 \text{ dynamics (from actuator)} \\
\dot{T}_2 &= T_2 \text{ dynamics (from actuator)}
\end{aligned}}
\end{equation}

\subsubsection{7.17.7 System Matrix [A] Derivation}

The attachment shows a system matrix $[A]$ for the linearized system. Let's derive this from the nonlinear dynamics.

\textbf{Nonlinear system:}

\begin{equation}
\dot{\mathbf{x}} = \mathbf{f}(\mathbf{x}, \mathbf{u})
\end{equation}

\textbf{Linearization around equilibrium point} $(\mathbf{x}_0, \mathbf{u}_0)$:

\begin{equation}
\delta \dot{\mathbf{x}} = \mathbf{A} \delta \mathbf{x} + \mathbf{B} \delta \mathbf{u}
\end{equation}

where:

\begin{equation}
\mathbf{A} = \frac{\partial \mathbf{f}}{\partial \mathbf{x}} \bigg|_{(\mathbf{x}_0, \mathbf{u}_0)}, \quad \mathbf{B} = \frac{\partial \mathbf{f}}{\partial \mathbf{u}} \bigg|_{(\mathbf{x}_0, \mathbf{u}_0)}
\end{equation}

\textbf{For a 6-state system} $\mathbf{x} = [x_m, y_m, \dot{x}_m, \dot{y}_m, T_1, T_2]^T$:

The Jacobian matrix $\mathbf{A}$ is $6 \times 6$:

\begin{equation}
\mathbf{A} = \frac{\partial \mathbf{f}}{\partial \mathbf{x}} = \begin{bmatrix}
\frac{\partial \dot{x}_m}{\partial x_m} & \frac{\partial \dot{x}_m}{\partial y_m} & \frac{\partial \dot{x}_m}{\partial \dot{x}_m} & \frac{\partial \dot{x}_m}{\partial \dot{y}_m} & \frac{\partial \dot{x}_m}{\partial T_1} & \frac{\partial \dot{x}_m}{\partial T_2} \\
\frac{\partial \dot{y}_m}{\partial x_m} & \frac{\partial \dot{y}_m}{\partial y_m} & \frac{\partial \dot{y}_m}{\partial \dot{x}_m} & \frac{\partial \dot{y}_m}{\partial \dot{y}_m} & \frac{\partial \dot{y}_m}{\partial T_1} & \frac{\partial \dot{y}_m}{\partial T_2} \\
\frac{\partial \ddot{x}_m}{\partial x_m} & \frac{\partial \ddot{x}_m}{\partial y_m} & \frac{\partial \ddot{x}_m}{\partial \dot{x}_m} & \frac{\partial \ddot{x}_m}{\partial \dot{y}_m} & \frac{\partial \ddot{x}_m}{\partial T_1} & \frac{\partial \ddot{x}_m}{\partial T_2} \\
\frac{\partial \ddot{y}_m}{\partial x_m} & \frac{\partial \ddot{y}_m}{\partial y_m} & \frac{\partial \ddot{y}_m}{\partial \dot{x}_m} & \frac{\partial \ddot{y}_m}{\partial \dot{y}_m} & \frac{\partial \ddot{y}_m}{\partial T_1} & \frac{\partial \ddot{y}_m}{\partial T_2} \\
\frac{\partial \dot{T}_1}{\partial x_m} & \frac{\partial \dot{T}_1}{\partial y_m} & \frac{\partial \dot{T}_1}{\partial \dot{x}_m} & \frac{\partial \dot{T}_1}{\partial \dot{y}_m} & \frac{\partial \dot{T}_1}{\partial T_1} & \frac{\partial \dot{T}_1}{\partial T_2} \\
\frac{\partial \dot{T}_2}{\partial x_m} & \frac{\partial \dot{T}_2}{\partial y_m} & \frac{\partial \dot{T}_2}{\partial \dot{x}_m} & \frac{\partial \dot{T}_2}{\partial \dot{y}_m} & \frac{\partial \dot{T}_2}{\partial T_1} & \frac{\partial \dot{T}_2}{\partial T_2}
\end{bmatrix}
\end{equation}

\textbf{Structured form matching the attachment:}

\begin{equation}
\mathbf{A} = \begin{bmatrix}
0 & 0 & 1 & 0 & 0 & 0 \\
0 & 0 & 0 & 1 & 0 & 0 \\
a_{31}(q_1,q_2) & a_{32}(q_1,q_2) & a_{33}(\dot{q}) & a_{34}(\dot{q}) & a_{35}(\mathbf{q}) & a_{36}(\mathbf{q}) \\
a_{41}(q_1,q_2) & a_{42}(q_1,q_2) & a_{43}(\dot{q}) & a_{44}(\dot{q}) & a_{45}(\mathbf{q}) & a_{46}(\mathbf{q}) \\
0 & 0 & 0 & 0 & -k_1 & 0 \\
0 & 0 & 0 & 0 & 0 & -k_2
\end{bmatrix}
\end{equation}

\subsubsection{7.17.8 Output Equation}

The attachment shows the output relationship. Typically, we measure end-effector position and velocity:

\begin{equation}
\mathbf{y} = \mathbf{C} \mathbf{x} + \mathbf{D} \mathbf{u}
\end{equation}

where:

\begin{equation}
\mathbf{C} = \begin{bmatrix}
1 & 0 & 0 & 0 & 0 & 0 \\
0 & 1 & 0 & 0 & 0 & 0 \\
0 & 0 & 1 & 0 & 0 & 0 \\
0 & 0 & 0 & 1 & 0 & 0 \\
0 & 0 & 0 & 0 & 1 & 0 \\
0 & 0 & 0 & 0 & 0 & 1
\end{bmatrix}
\end{equation}

(Identity matrix - we measure all states)

\begin{equation}
\mathbf{D} = \mathbf{0}_{6 \times 2}
\end{equation}

(No direct feedthrough)

\subsubsection{7.17.9 Complete Linear State-Space System}

\textbf{Standard form:}

\begin{equation}
\boxed{\dot{\mathbf{x}} = \mathbf{A} \mathbf{x} + \mathbf{B} \mathbf{u}}
\end{equation}

\begin{equation}
\boxed{\mathbf{y} = \mathbf{C} \mathbf{x} + \mathbf{D} \mathbf{u}}
\end{equation}

where:
\begin{itemize}
    \item $\mathbf{x} = [x_m, y_m, \dot{x}_m, \dot{y}_m, T_1, T_2]^T$ (6-state vector)
    \item $\mathbf{u} = [\tau_1, \tau_2]^T$ or $[\tau_{c,1}, \tau_{c,2}]^T$ (2-input vector)
    \item $\mathbf{y} = [x_m, y_m, \dot{x}_m, \dot{y}_m, T_1, T_2]^T$ (6-output vector)
    \item $\mathbf{A}$ is $6 \times 6$ (system dynamics matrix)
    \item $\mathbf{B}$ is $6 \times 2$ (input matrix)
    \item $\mathbf{C}$ is $6 \times 6$ (output matrix)
    \item $\mathbf{D}$ is $6 \times 2$ (feedthrough matrix)
\end{itemize}

\subsubsection{7.17.10 Specific Components of Matrices for Planar 2-DOF Arm}

\textbf{System matrix elements:}

From the dynamics, the upper block rows (position and velocity equations) give:

\begin{equation}
\begin{bmatrix}
\dot{x}_m \\
\dot{y}_m \\
\dot{x}_{m,1} \\
\dot{y}_{m,1}
\end{bmatrix} = \begin{bmatrix}
0 & 0 & 1 & 0 & 0 & 0 \\
0 & 0 & 0 & 1 & 0 & 0 \\
a_{31} & a_{32} & a_{33} & a_{34} & a_{35} & a_{36} \\
a_{41} & a_{42} & a_{43} & a_{44} & a_{45} & a_{46}
\end{bmatrix} \begin{bmatrix}
x_m \\
y_m \\
\dot{x}_m \\
\dot{y}_m \\
T_1 \\
T_2
\end{bmatrix} + \mathbf{B}_{\text{upper}} \mathbf{u}
\end{equation}

For the tension rows (assuming first-order dynamics):

\begin{equation}
\begin{bmatrix}
\dot{T}_1 \\
\dot{T}_2
\end{bmatrix} = \begin{bmatrix}
0 & 0 & 0 & 0 & -k_1 & 0 \\
0 & 0 & 0 & 0 & 0 & -k_2
\end{bmatrix} \begin{bmatrix}
x_m \\
y_m \\
\dot{x}_m \\
\dot{y}_m \\
T_1 \\
T_2
\end{bmatrix} + \mathbf{B}_{\text{lower}} \mathbf{u}
\end{equation}

\subsubsection{7.17.11 Comparison with Attachment Equations}

Your attachment shows the state vector organized as:

\begin{equation}
[A] \begin{bmatrix}
x_m \\ y_m \\ x_{m,1} \\ y_{m,1} \\ T_1 \\ T_2
\end{bmatrix} = \begin{bmatrix}
0 \\ m_1 g \\ \frac{1}{2} m_1 g \cos q_1 \\ 0 \\ m_2 g \\ \frac{1}{2} m_2 g \cos(q_1 + q_2)
\end{bmatrix}
\end{equation}

This can be interpreted as:

\textbf{Equation 1:} $0 = \dot{x}_m$ (equilibrium condition for position)

\textbf{Equation 2:} $m_1 g = 0$ (gravitational term for link 1)

And so on for the other equations. This suggests your attachment is showing an **equilibrium analysis** or **control input decomposition**, where the RHS shows the gravity and inertial effects.

\subsubsection{7.17.12 Connecting to Your Attachment}

If we reorganize the Euler-Lagrange equations to separate control inputs from dynamics:

\begin{equation}
\mathbf{M}(\mathbf{q}) \ddot{\mathbf{q}} = \boldsymbol{\tau} + \mathbf{J}_{\rho}^T \mathbf{T} - \mathbf{G}(\mathbf{q})
\end{equation}

This can be rewritten in control form as:

\begin{equation}
\ddot{\mathbf{q}} = \mathbf{M}^{-1}(\mathbf{q}) [\text{Actuator Forces}] - \mathbf{M}^{-1}(\mathbf{q}) \mathbf{G}(\mathbf{q})
\end{equation}

where the actuator forces are the cable tension contributions:

\begin{equation}
\text{Actuator Forces} = \mathbf{J}_{\rho}^T(\mathbf{q}) \mathbf{T} + \boldsymbol{\tau}
\end{equation}

Your attachment's matrix $[A]$ appears to decompose this into individual components:

\begin{itemize}
    \item Row 1: Position dynamics for $x_m$
    \item Row 2: Gravitational effect for link 1
    \item Row 3: Acceleration component from inertia
    \item Row 4: Gravity component for position y
    \item Row 5: Gravity for link 2
    \item Row 6: Acceleration from combined inertia
\end{itemize}

The specific gravity terms $(\frac{1}{2}m_i g \cos q_i)$ come directly from $\mathbf{G}(\mathbf{q})$ as derived in Section 7.8.3.

\subsubsection{7.17.13 Numerical Implementation}

To implement these equations programmatically:

\textbf{Python pseudocode:}

\begin{lstlisting}[language=Python]
import numpy as np

class HybridRobotDynamics:
    def __init__(self, l1, l2, m1, m2, I1, I2, g=9.81):
        self.l1, self.l2 = l1, l2
        self.m1, self.m2 = m1, m2
        self.I1, self.I2 = I1, I2
        self.g = g
    
    def forward_kinematics(self, q):
        x = self.l1*np.cos(q[0]) + self.l2*np.cos(q[0]+q[1])
        y = self.l1*np.sin(q[0]) + self.l2*np.sin(q[0]+q[1])
        return np.array([x, y])
    
    def velocity_kinematics(self, q, dq):
        # Jacobian matrix
        J = np.array([
            [-self.l1*np.sin(q[0]) - self.l2*np.sin(q[0]+q[1]), 
             -self.l2*np.sin(q[0]+q[1])],
            [self.l1*np.cos(q[0]) + self.l2*np.cos(q[0]+q[1]),
             self.l2*np.cos(q[0]+q[1])]
        ])
        return J @ dq
    
    def mass_matrix(self):
        M = np.array([
            [self.I1 + self.I2, self.I2],
            [self.I2, self.I2]
        ])
        return M
    
    def gravity_torques(self, q):
        g1 = (self.m1*self.l1/2 + self.m2*self.l1)*self.g*np.cos(q[0]) \
           + self.m2*self.l2/2*self.g*np.cos(q[0]+q[1])
        g2 = self.m2*self.l2/2*self.g*np.cos(q[0]+q[1])
        return np.array([g1, g2])
    
    def state_derivative(self, x, u, T):
        # x = [x_m, y_m, dx_m, dy_m, T1, T2]
        # u = [tau1, tau2]
        q = self.compute_q_from_xm(x[:2])
        dq = self.compute_dq_from_dxm(q, x[2:4])
        
        M = self.mass_matrix()
        G = self.gravity_torques(q)
        J_rho_T = self.cable_geometry_jacobian_transpose(q, T)
        
        # Solve for accelerations
        ddq = np.linalg.solve(M, u + J_rho_T - G)
        
        # Compute accelerations in Cartesian
        dx_m = x[2:4]
        ddx_m = self.compute_cartesian_acceleration(q, dq, ddq)
        
        # Tension dynamics
        dT = self.tension_dynamics(T, u_tension)
        
        return np.concatenate([dx_m, ddx_m, dT])
\end{lstlisting}

\subsubsection{7.17.14 Connection to Singularity Analysis}

Near singularities, the system exhibits critical behavior:

\begin{enumerate}
    \item \textbf{Rank deficiency in Jacobian}: $\text{rank}(\mathbf{J}_{\rho}) < 2$
    
    \item \textbf{Large accelerations from small inputs}: 
    \begin{equation}
    \|\ddot{\mathbf{q}}\| = \left\|\mathbf{M}^{-1} \mathbf{J}_{\rho}^T \mathbf{T}\right\| \to \infty \text{ as } \det(\mathbf{J}_{\rho}) \to 0
    \end{equation}
    
    \item \textbf{Unstable pole near singularity}: The linearized system matrix $\mathbf{A}$ has eigenvalues with large magnitude
    
    \item \textbf{Loss of controllability}: System cannot be controlled in all directions
    
    \begin{equation}
    \text{rank}[\mathbf{B}, \mathbf{A}\mathbf{B}, \mathbf{A}^2\mathbf{B}, \ldots] < n \text{ at singularity}
    \end{equation}
\end{enumerate}

This reinforces the importance of singularity detection and avoidance algorithms.

\subsubsection{7.17.15 Summary: From Euler-Lagrange to State-Space}

The derivation path is:

\begin{enumerate}
    \item **Define state vector**: $\mathbf{x} = [x_m, y_m, \dot{x}_m, \dot{y}_m, T_1, T_2]^T$
    
    \item **Write Lagrangian**: $L = T - V$ from kinetic and potential energies
    
    \item **Apply Euler-Lagrange**: $\frac{d}{dt}\frac{\partial L}{\partial \dot{\mathbf{q}}} - \frac{\partial L}{\partial \mathbf{q}} = \mathbf{Q}$
    
    \item **Get dynamics equations**: $\mathbf{M}(\mathbf{q})\ddot{\mathbf{q}} + \mathbf{G}(\mathbf{q}) = \boldsymbol{\tau} + \mathbf{J}_{\rho}^T\mathbf{T}$
    
    \item **Transform to Cartesian**: Convert joint accelerations to end-effector accelerations via kinematics
    
    \item **Add tension dynamics**: Include first-order models for cable tension control
    
    \item **Linearize for analysis**: Compute $\mathbf{A}$ matrix for stability and control analysis
    
    \item **Result**: Complete state-space system for simulation and control design
\end{enumerate}
